\documentclass[12pt,a4paper]{article}
\newcommand{\StudentName}{Nguyen Thai Thanh Binh}
\newcommand{\StudentID}{2570175}
\newcommand{\ReportVariant}{Analytical derivation}
\usepackage{ee5430_style}

\begin{document}
\MakeTitlePage
\tableofcontents
\newpage

\section{Problem 1: Parameters of a lossy transmission line}

\LineElementDiagram

The distributed parameters are
\[
R=1.5~\Omega/\mathrm{m},\quad G=0.4~\mathrm{mS/m},\quad
L=10~\mathrm{nH/m},\quad C=0.2~\mathrm{pF/m},\quad f=2~\mathrm{GHz}.
\]
With the $\ee^{\jmath\omega t}$ convention, the exact characteristic impedance and
propagation constant are
\begin{align}
Z_0&=\sqrt{\frac{R+\jmath\omega L}{G+\jmath\omega C}},\\
\gamma&=\alpha+\jmath\beta
=\sqrt{(R+\jmath\omega L)(G+\jmath\omega C)}.
\end{align}
At $\omega=2\pi f=1.2566371\times10^{10}~\mathrm{rad/s}$,
\begin{align*}
R+\jmath\omega L&=1.5+\jmath125.6637~\Omega/\mathrm{m},\\
G+\jmath\omega C&=4.0000\times10^{-4}
 +\jmath2.51327\times10^{-3}~\mathrm{S/m}.
\end{align*}
Therefore,
\begin{align}
Z_0&=221.6297+\jmath16.1961~\Omega
     =222.2207\angle4.1796\degree~\Omega,\\
\gamma&=0.0479467+\jmath0.5634947~\mathrm{m}^{-1}.
\end{align}
Consequently,
\begin{align*}
\alpha&=0.0479467~\mathrm{Np/m},\\
\alpha_{\mathrm{dB}}&=8.6858896\alpha=0.41646~\mathrm{dB/m},\\
\beta&=0.563495~\mathrm{rad/m},\\
\lambda&=\frac{2\pi}{\beta}=11.1504~\mathrm{m}.
\end{align*}

\answer{$Z_0=221.63+\jmath16.20~\Omega$ and
$\gamma=(0.04795+\jmath0.56349)~\mathrm{m}^{-1}$.  Thus the line attenuates
at $0.04795~\mathrm{Np/m}=0.4165~\mathrm{dB/m}$, while
$\beta=0.56349~\mathrm{rad/m}$ and $\lambda=11.150~\mathrm{m}$.}

\section{Problem 2: A nominal 50-ohm coaxial cable at 500 MHz}

Here
\[
R=4~\Omega/\mathrm{m},\quad G=1~\mathrm{mS/m},\quad
L=250~\mathrm{nH/m},\quad C=100~\mathrm{pF/m},\quad f=500~\mathrm{MHz}.
\]

\subsection{Exact calculation}

For $\omega=3.1415927\times10^9~\mathrm{rad/s}$,
\begin{align*}
R+\jmath\omega L&=4+\jmath785.3982~\Omega/\mathrm{m},\\
G+\jmath\omega C&=0.001+\jmath0.3141593~\mathrm{S/m}.
\end{align*}
Substitution in the exact formulas gives
\begin{align}
Z_0&=50.0001748-\jmath0.0477458~\Omega
    =50.0001976\angle(-0.054712\degree)~\Omega,\\
\gamma&=0.06499997+\jmath15.70797043~\mathrm{m}^{-1}.
\end{align}
Thus $\alpha=0.06499997~\mathrm{Np/m}=0.5645897~\mathrm{dB/m}$ and
$\beta=15.70797043~\mathrm{rad/m}$.

\subsection{Low-loss approximation and comparison}

The loss ratios are
\[
\frac{R}{\omega L}=5.093\times10^{-3}\ll1,
\qquad
\frac{G}{\omega C}=3.183\times10^{-3}\ll1,
\]
so the low-loss formulas apply:
\begin{align}
Z_0&\simeq\sqrt{\frac{L}{C}}=50.000~\Omega,\\
\alpha&\simeq\frac{1}{2}\left(\frac{R}{Z_0}+GZ_0\right)
=\frac{1}{2}(0.08+0.05)=0.065000~\mathrm{Np/m},\\
\beta&\simeq\omega\sqrt{LC}=15.7079633~\mathrm{rad/m}.
\end{align}

\begin{center}
\renewcommand{\arraystretch}{1.2}
\begin{tabular}{lccc}
\toprule
Quantity & Exact & Low-loss & Relative difference \\
\midrule
$|Z_0|$ ($\Omega$) & 50.0001976 & 50.0000000 & $3.95\times10^{-4}\%$ \\
$\alpha$ (Np/m) & 0.06499997 & 0.06500000 & $4.56\times10^{-5}\%$ \\
$\beta$ (rad/m) & 15.70797043 & 15.70796327 & $4.56\times10^{-5}\%$ \\
\bottomrule
\end{tabular}
\end{center}

\subsection{Phase velocity and wavelength}

Using the exact phase constant,
\begin{align}
v_p&=\frac{\omega}{\beta}=1.9999991\times10^8~\mathrm{m/s},\\
\lambda&=\frac{2\pi}{\beta}=0.3999998~\mathrm{m}.
\end{align}
The low-loss values are $v_p\simeq1/\sqrt{LC}=2.0000\times10^8~\mathrm{m/s}$
and $\lambda\simeq0.4000~\mathrm{m}$.

\answer{Exact: $Z_0=50.00017-\jmath0.04775~\Omega$ and
$\gamma=(0.06499997+\jmath15.70797043)~\mathrm{m}^{-1}$.  The low-loss
approximation is essentially identical; $v_p\approx2.00\times10^8~\mathrm{m/s}$
and $\lambda\approx0.400~\mathrm{m}$.}

\section{Problem 3: Matched, lossless line}

\MatchedLineDiagram{v_s(t)=5\cos(2\pi\,2\!\times\!10^8t)}{75~\Omega}{50~\Omega}{3~\mathrm{m}}{R_L=50~\Omega}

The load is matched, so $\Gamma_L=0$ and only a forward wave exists.  The
source frequency and the line constants are
\begin{align*}
f&=2.0\times10^8~\mathrm{Hz}, &
\omega&=4\pi\times10^8~\mathrm{rad/s},\\
v_p&=2.0\times10^8~\mathrm{m/s}, &
\lambda&=\frac{v_p}{f}=1.0~\mathrm{m},\\
\beta&=\frac{2\pi}{\lambda}=2\pi~\mathrm{rad/m}.&&
\end{align*}
Because the input impedance of a matched line equals $Z_0$, the incident-wave
amplitude follows from a voltage divider:
\[
V_0^+=5\frac{Z_0}{R_s+Z_0}
=5\frac{50}{75+50}=2.0~\mathrm{V},
\qquad I_0^+=\frac{V_0^+}{Z_0}=0.040~\mathrm{A}.
\]
Taking $z=0$ at the line input, the instantaneous waves are therefore
\begin{align}
v(z,t)&=2\cos\left(4\pi\times10^8t-2\pi z\right)~\mathrm{V},\\
i(z,t)&=0.040\cos\left(4\pi\times10^8t-2\pi z\right)~\mathrm{A},
\qquad 0\le z\le3~\mathrm{m}.
\end{align}
Since $\ell/\lambda=3$, the electrical length is
$\beta\ell=6\pi~\mathrm{rad}=1080\degree$.  At the load,
\[
v_L(t)=v(3,t)=2\cos(4\pi\times10^8t-6\pi)
=2\cos(4\pi\times10^8t)~\mathrm{V}.
\]

\answer{The line is three wavelengths long.  Its load voltage has a peak
value of $2~\mathrm{V}$ and is in phase with the source voltage (the accumulated
phase is $1080\degree$, an integer number of cycles).}

\section{Problem 4: Matched air line and an open-circuit comparison}

\MatchedLineDiagram{v_s(t)=10\cos(2\pi\,5\!\times\!10^7t)}{25~\Omega}{100~\Omega}{6~\mathrm{m}}{R_L=100~\Omega}

For an air line, $v_p=3.0\times10^8~\mathrm{m/s}$.  Hence
\begin{align*}
f&=5.0\times10^7~\mathrm{Hz}, & \omega&=\pi\times10^8~\mathrm{rad/s},\\
\lambda&=\frac{v_p}{f}=6.0~\mathrm{m}, &
\beta&=\frac{2\pi}{\lambda}=\frac{\pi}{3}~\mathrm{rad/m}.
\end{align*}

\subsection{Instantaneous voltage and current}

For the matched load, $Z_{\mathrm{in}}=Z_0=100~\Omega$ and
\[
V_0^+=10\frac{100}{25+100}=8.0~\mathrm{V},
\qquad I_0^+=\frac{8}{100}=0.080~\mathrm{A}.
\]
Thus
\begin{align}
v(z,t)&=8\cos\left(\pi\times10^8t-\frac{\pi}{3}z\right)~\mathrm{V},\\
i(z,t)&=0.080\cos\left(\pi\times10^8t-\frac{\pi}{3}z\right)~\mathrm{A},
\qquad 0\le z\le6~\mathrm{m}.
\end{align}

\subsection{Average load power}

The load voltage has peak amplitude $8~\mathrm{V}$, so
\begin{equation}
P_{L,\mathrm{av}}=\frac{V_{L,\mathrm{rms}}^2}{R_L}
=\frac{(8/\sqrt2)^2}{100}=0.320~\mathrm{W}.
\end{equation}

\subsection{If the load is replaced by an open circuit}

The forward-wave expression above is no longer valid because an open circuit
has $\Gamma_L=+1$.  A reflected wave of equal voltage amplitude is produced.
Here $\ell=\lambda$, so $Z_{\mathrm{in}}\to\infty$ and the source current is
zero; therefore the total input voltage equals the source voltage.  With
$\phasor V_0^-=\phasor V_0^+=5~\mathrm{V}$,
\begin{align}
\phasor V(z)&=5\ee^{-\jmath\beta z}+5\ee^{\jmath\beta z}
=10\cos(\beta z),\\
v(z,t)&=10\cos\left(\frac{\pi z}{3}\right)
\cos(\pi\times10^8t)~\mathrm{V}.
\end{align}
This is a standing wave, not a single traveling wave.

\answer{For the matched load, $v(z,t)=8\cos(\pi10^8t-\pi z/3)$ V,
$i(z,t)=0.08\cos(\pi10^8t-\pi z/3)$ A, and $P_L=0.320$ W.
With an open circuit the formula changes to a standing-wave expression because
$\Gamma_L=+1$.}

\section*{Final numerical summary}
\addcontentsline{toc}{section}{Final numerical summary}
\begin{center}
\renewcommand{\arraystretch}{1.25}
\begin{tabular}{>{\bfseries}c p{0.77\linewidth}}
\toprule
Problem & Main result \\
\midrule
1 & $Z_0=221.63+\jmath16.20~\Omega$;
$\gamma=0.04795+\jmath0.56349~\mathrm{m}^{-1}$;
$\lambda=11.150~\mathrm{m}$. \\
2 & $Z_0=50.00017-\jmath0.04775~\Omega$;
$\gamma=0.06500+\jmath15.70797~\mathrm{m}^{-1}$;
$v_p\approx2.00\times10^8~\mathrm{m/s}$, $\lambda\approx0.400~\mathrm{m}$. \\
3 & $v=2\cos(4\pi10^8t-2\pi z)$ V,
$i=0.04\cos(4\pi10^8t-2\pi z)$ A; $\ell=3\lambda$. \\
4 & $v=8\cos(\pi10^8t-\pi z/3)$ V,
$i=0.08\cos(\pi10^8t-\pi z/3)$ A; $P_L=0.320$ W. \\
\bottomrule
\end{tabular}
\end{center}

\end{document}
